% Generated draft: requires mathematical and source review.
\documentclass[11pt]{article}
\usepackage[margin=25mm]{geometry}
\usepackage{fontspec}
\setmainfont{Latin Modern Roman}
\usepackage{amsmath,amssymb,mathtools,unicode-math}
\setmathfont{Latin Modern Math}
\usepackage{enumitem,needspace,xurl,hyperref}
\setlength{\parindent}{0pt}
\setlength{\parskip}{5pt}
\setlength{\emergencystretch}{3em}
\newcommand{\N}{\mathbb{N}}
\newcommand{\Z}{\mathbb{Z}}
\newcommand{\Q}{\mathbb{Q}}
\newcommand{\R}{\mathbb{R}}
\newcommand{\C}{\mathbb{C}}
\newcommand{\F}{\mathbb{F}}
\newcommand{\A}{\mathbb{A}}
\newcommand{\PP}{\mathbb{P}}
\newcommand{\E}{\mathbb{E}}
\newcommand{\eps}{\varepsilon}
\newcommand{\dd}{\,\mathrm{d}}
\begin{document}
\section*{UnsolvedMath ID: 30004325}
\textbf{Source problem code:} catalogue entry
\subsection{Definitions and mathematical statement}
Let \(X\) be a smooth projective surface obtained by blowing up \(\PP^2\) at finitely many very general points, and let \(L\) be an ample line bundle on \(X\). The source specifies neither the base field nor the number of points. In the usual geometric interpretation, a very general configuration lies outside a countable union of proper Zariski closed subsets of the parameter space of configurations; the base field and the existence of such configurations require clarification.

For a point \(x\in X\), the usual pointwise Seshadri constant is
\[
\varepsilon(L;x)=\inf_{\substack{C\subset X\text{ reduced irreducible curve}\\x\in C}}\frac{L\cdot C}{\operatorname{mult}_x(C)},
\]
where \(L\cdot C\) is the intersection number and \(\operatorname{mult}_x(C)\) is the multiplicity of \(C\) at \(x\). Under the usual global convention,
\[
\varepsilon(X,L)=\inf_{x\in X}\varepsilon(L;x).
\]
Here the points and curves are understood in the intended geometric setting, which the supplied record does not specify. In particular, no base field or base-change convention is imposed by this draft.

With this global convention, the supplied question asks whether
\[
\varepsilon(X,L)\in\Q
\quad\text{for every such pair }(X,L),
\]
or whether there exists such a pair with \(\varepsilon(X,L)\notin\Q\). The source does not explicitly define ``global,'' so its intended meaning must be checked against the report. The infimum over all points must be distinguished from the value at a very general point. The supplied literature discussion concerns conditional irrationality implications from strengthened SHGH or Nagata conjectures; it does not specify those strengthened hypotheses or explain their precise relation to the global formulation.
\subsection{Short English statement}
For a smooth projective plane blowup at very general points and an ample line bundle, must the global Seshadri constant be rational, or can it be irrational? This draft uses the usual infimum over all points; the intended convention and base field require review.
\subsection{Sources}
\begin{itemize}
\item Supplied local UnsolvedMath registry record, permanent ID \(30004325\), source problem code \texttt{catalogue entry}, title ``Rationality of Seshadri Constants on Plane Blowups,'' in the Oberwolfach Reports Open Problems collection. No individual proposer is supplied.
\item Supplied report citation: \emph{Mini-Workshop: Seshadri Constants (2020)}, Oberwolfach Reports, pp.\ 3271--3313, DOI \texttt{10.4171/owr/2019/53}. \url{https://doi.org/10.4171/owr/2019/53}. The record gives proposed year 2019, while the citation labels the workshop 2020; this discrepancy is unresolved.
\item Supplied literature reference: ``Single-point Seshadri constants on rational surfaces.'' \url{https://digitalcommons.unl.edu/mathfacpub/120/}.
\item Supplied literature reference: ``Seshadri constants and related conjectures over characteristic zero.'' \url{https://www.sciencedirect.com/science/article/pii/S0007449725000430}.
\item Supplied literature reference: ``Oberwolfach Reports, irrationality discussion.'' \url{https://doi.org/10.4171/owr/2019/53}.
\item Literature assessment dated exactly \texttt{2026-08-22}: the supplied curation labels the problem open and states that the cited implications give irrational Seshadri constants only conditionally on strengthened SHGH or Nagata conjectures, supplying no unconditional example and leaving the general irrationality question open. This is supplied dated curation derived from the OWR material, not an independently verified assessment or verification of the precise global formulation. The record expressly requests independent verification of time-sensitive status and cited claims.
\item The supplied statement status is \texttt{corrected\_verified}; its accompanying qualification says that only formula delimiters, notation, grammar, and typography were repaired, with mathematical content and scope preserved. This label does not establish mathematical correctness.
\end{itemize}
\end{document}
